% Appendix N: Learning Geometry and Curriculum Dynamics

\section*{N.1 Introduction}

Traditional educational theory assumes knowledge consists of facts
accumulated over time. The accessibility framework suggests otherwise:
\emph{knowledge is not accumulated — knowledge is navigated.
Education becomes the engineering of trajectories through semantic space.}

\begin{definition}{Educational State Space}{edu-state}
  A learner occupies $x(t) \in M$ (semantic manifold).
  Learning is motion through $M$; the learner is a trajectory.
\end{definition}

\begin{definition}{Accessibility Value of a Concept}{concept-access}
  $S(c) = \log |\mathcal{A}(c)|$, where $\mathcal{A}(c)$ is the set of
  future concepts reachable from $c$. Concepts with large accessibility
  act as gateways (algebra, geometry, logic, probability, programming).
\end{definition}

\begin{definition}{Curriculum as Vector Field}{curriculum-field}
  A \emph{curriculum field} $\mathbf{C}(x)$ guides educational motion.
  Different curricula correspond to different vector fields; the same
  learner may reach different destinations depending on the field.
\end{definition}

\begin{example}
  Geometry-first: $\text{geometry} \to \text{trigonometry} \to
  \text{calculus} \to \text{differential geometry}$.\\
  Algebra-first: $\text{algebra} \to \text{functions} \to
  \text{calculus} \to \text{abstract algebra}$.\\
  Logic-first: favors formal reasoning, theorem construction,
  computational thinking.\\
  Statistics-first: probabilistic rather than deterministic reasoning.
  Each produces a distinct accessibility landscape.
\end{example}

\begin{definition}{Educational Phase Transition}{edu-phase}
  A \emph{phase transition} occurs when a small conceptual change produces
  a large increase in accessibility (learning algebra, understanding proof,
  mastering recursion, discovering coordinate systems).
\end{definition}

\begin{definition}{Optimal Curriculum}{opt-curriculum}
  An \emph{optimal curriculum} maximizes
  $\int_0^T S_E(t)\,dt$,
  where $S_E = \log |\mathcal{A}(x)|$ is educational entropy.
  The objective is not memorization; it is future opportunity.
\end{definition}

\begin{theorem}{Curriculum Divergence Theorem}{curriculum-div}
  Distinct curriculum fields generally produce distinct cognitive
  trajectories.
\end{theorem}
\begin{proof}
  Distinct vector fields $\mathbf{C}_1 \neq \mathbf{C}_2$ generate
  distinct integral curves except in degenerate cases.
\end{proof}

\begin{namedthm}{Educational Accessibility Principle}
  The purpose of education is not information transfer. The purpose of
  education is accessibility expansion:
  \[
    \text{Education} \to \text{Conceptual Navigation}
    \to \text{Accessibility Expansion} \to \text{Agency.}
  \]
\end{namedthm}
